\chapter{Nut Allergies in Children}
\index{Nut Allergies in Children}

\medskip
\noindent\textit{\textbf{Under review.} This chapter is an AI-assisted educational draft; it is not a substitute for professional medical advice.}

\section*{Definition and Overview}
Nut allergy is an immune reaction to one or more tree nuts or peanuts. Reactions can be mild or can progress rapidly to anaphylaxis. A suspected allergy should be assessed by a qualified clinician; families should not deliberately re-challenge a child at home.

\section*{Clinical Features}
Symptoms may include hives, swelling of the lips or face, vomiting, cough, wheeze, hoarse voice, breathing difficulty, dizziness, or collapse. Anaphylaxis can occur without skin symptoms.

\section*{Emergency Management}
If anaphylaxis is suspected, give the child's prescribed adrenaline injector immediately, contact local emergency services, lay the child flat unless breathing is difficult, and follow the injector action plan. Give a second injector if advised and symptoms persist while waiting for help.

\section*{Diagnosis}
Assessment includes the reaction history, examination, food and cofactor history, and selected skin-prick or blood tests. Oral food challenge is performed only under specialist supervision.

\section*{Management}
Avoid the confirmed trigger while maintaining a nutritionally adequate diet. Provide an ASCIA action plan, prescribe and teach use of adrenaline injectors when indicated, inform school or childcare, and review the plan regularly. Do not rely on antihistamines to treat anaphylaxis.

\section*{Prevention and Self-Care}
Read labels, prevent cross-contact, teach carers and the child age-appropriate safety, and seek review after any reaction. Emergency advice takes priority over this general information.
\section*{Safety-netting}
Seek urgent care for breathing difficulty, blue or very pale appearance, seizure, reduced consciousness, severe dehydration, serious injury, or safeguarding concerns.
